\section{Deployment, Calibration and Explanation Studies}
\label{sec:phase2}
The experiments in this section were added after the main evaluation. Each has
a protocol committed to the repository before its first result
(\texttt{docs/research/}), and every amendment is recorded there with its
reason.

\subsection{RQ6: latency on the deployed stack}
The latencies of Section~\ref{sec:results} come from the in-memory store. We
therefore replayed the paced scenarios of RQ1 and the users sweep through
the deployed path: Redis Streams into the detection engine, with the graph in
Neo4j~5.26. All \EightRuns{} runs produced the same sequence of ring updates and
the same flagged rings as the in-memory engine (\EightSame{} of \EightRuns{}),
so the store and the bus change timing, never results. Paced at simulated real
time, a flag followed the event that completed its batch by a median of
\EightFlagLo{}--\EightFlagHi{}\,ms across the five scenarios. Each detection
reads its region through Cypher round trips and took \EightSlowLo{} to
\EightSlowHi{} times as long as on the in-memory store for the same batches.
When events arrive faster than the store can take them, Neo4j sets the ceiling:
\EightTputLo{}--\EightTputHi{} events per second on the test laptop, with events
queueing in Redis and flags arriving \EightBurstLo{}--\EightBurstHi{}\,s after
their event. Varying the consumer's read size and blocking timeout, after
HYADES~\cite{bora2026hyades}, changed nothing measurable at these rates. Neo4j
saturates before the consumer does.

\subsection{RQ7: ring-level calibration and early warning}
The scorer's isotonic calibration was fitted on snapshot rows, but analysts read
scores of ring-level alerts. On \TenUnits{} ring-level units from the test seeds
its expected calibration error is \TenEce{}, and rings scored in the 0.8--0.9
and 0.9--1.0 bins were true rings \TenBinEight\% and \TenBinNine\% of the time.
A calibrator refitted on ring-level units of separate seeds, in the manner
compared by Azbergen et al.~\cite{azbergen2026calibration} and Ding et
al.~\cite{ding2021temperature}, lowers the Brier score from \TenBrier{} to
\TenBrierRefit{}. Ranking holds up better than calibration: precision at the top
5\% of units is \TenPAtFive{}. Early warning, measured without look-ahead, is
rare. Of \ElevenEligible{} true rings with a forming phase, \ElevenWarned{}
(\ElevenShare\%) were flagged before activation, a median of \ElevenLead{}\,s
early, and \ElevenPrecision\% of first early warnings pointed at a real ring.

\subsection{RQ8: explanation quality}
We measured the explanations the product shows on the \TwelveAlerts{} alerts of
the test streams. The unit is an alert (a ring at its peak calibrated score of
at least \TunedAlert{}); a model counterfactual is a proposed change of one
actionable feature, and an alert can have none, one or several.
\TwelveCfAlerts{} of the \TwelveAlerts{} alerts have at least one, \TwelveCfs{}
in all. Writing the unit test for counterfactual validity exposed a defect:
counterfactuals aimed at the raw score of the calibrator's breakpoint, but the
calibrator interpolates between breakpoints, so the proposed change left the
calibrated score above the threshold. As shipped, \TwelveValidAsShipped{} of
the \TwelveCfs{} counterfactuals flipped the decision; aimed at the
interpolated crossing, \TwelveValidCorrected{} of \TwelveCfs{} do. The
correction was registered as a protocol amendment before the experiment ran.
Any system that derives counterfactual targets from an isotonic calibrator can
make the same mistake.

Coverage is the larger limitation: the other \TwelveNoCfShare\% of alerts have
no single-feature counterfactual over the actionable features, because typical
transfer size dominates their score. Under 5\% noise in non-actionable
features, after Artelt et al.~\cite{artelt2021robustness}, the contribution
vector moves by a median of \TwelveDriftFive{} (relative $L_1$) and the top
three features never change for the median alert; for a linear model this is
close to guaranteed, so it mainly checks the implementation. A scorer retrained
on disjoint seeds, after Movin et al.~\cite{movin2025acdc}, ranks the
contributions similarly (median Spearman \TwelveSpearman{}) and still raises
\TwelveStillAlert\% of the alerts, but it proposes the same first
counterfactual for only \TwelveCfSame\% of them, so a counterfactual describes
this model, not the ring. Against simulator ground truth, the members shown in
the evidence subgraph of a correct alert are the true ring's (median precision
\TwelveFidPrec{}, coverage \TwelveFidCov{}). Because that subgraph is built from
the flagged ring's members, this checks that the evidence shows the right
accounts; it does not test whether the evidence explains the score. A
deletion test would: remove the shared-attribute links or the internal
transfers that the evidence lists, recompute the ring's features and score, and
compare the drop with removing as many links of the same members that the
evidence does not list. We have not run it.

\subsection{RQ9: a learned baseline and evasion strength}
To test whether the rule stage is needed, we trained gradient boosting on the
same 19 ring features over rule-free candidate groups (Louvain communities
without a density filter). Training, calibration and threshold selection used
the shipped scorer's protocol. On the test streams of RQ1 it reaches F1
\ThirteenGbF{} against \ThirteenPropF{} (paired Wilcoxon $p\ThirteenP$,
rank-biserial \ThirteenRb{}). Precision is not the problem: it is \ThirteenGbP{}
against \ThirteenPropP{}, with a third of the alerts. Recall is:
\ThirteenGbR{} against \ThirteenPropR{}. The gap is largest where hubs and
evasion matter (\texttt{hub\_heavy}: \ThirteenGbFHub{} against
\ThirteenPropFHub{}; \texttt{adversarial}: \ThirteenGbFAdv{} against
\ThirteenPropFAdv{}).

Turning up camouflage strength, after Liu et al.~\cite{liu2025robust} and Hu et
al.~\cite{hu2025hhgfd}, separates the methods further. With up to eight decoy
edges per ring action, the proposed method's recall of the decoy rings stays
between \FourteenDecoyPropLo{} and \FourteenDecoyPropHi{}, while Louvain falls
to \FourteenDecoyLouvainMax{} and the learned baseline to
\FourteenDecoyGbMax{}. Slow rings remain the weakness found in RQ3: the proposed
method recalls at most \FourteenSlowPropHi{} of them, and static Cypher with a
degree cap of ten recovers more (\FourteenSlowCapLo{}--\FourteenSlowCapHi{}),
though at a precision near 0.04. In the 60\,s scenario, strong slow and
identity-rotation settings produce few rings, so those two sweeps have little
statistical power. We report them as indicative.
